\documentclass[11pt,a4paper]{article}
\usepackage{fontspec}
\setmainfont{Latin Modern Roman}
\usepackage[margin=2.2cm]{geometry}
\usepackage{booktabs,enumitem}
\usepackage[hidelinks]{hyperref}
\setlist{itemsep=1pt,topsep=2pt,leftmargin=*}
\setlength{\parskip}{3pt}
\setlength{\parindent}{0pt}

\begin{document}
\begin{center}
{\Large\bfseries ResearchLens: Recommending What to Read Next\\from a Library of Saved Papers}\par\medskip
Progress report, IT4868E Web Mining, semester 2026.1. Group 5\par
Tran Quang Hung (20235502), Vu Thuong Tin (20230091),\\ Chu Anh Duc (20230081), Dinh Quoc Bao (20235478)\par
Instructor: Dr. Nguyễn Kiêm Hiếu. Hanoi, 23 November 2026
\end{center}

\section{Problem statement}
Given a few papers a researcher has saved, recommend which paper to read next, and show why: papers that cite it together with a saved paper, its citation links to the library, and its peer reviews. The project covers link analysis, recommender systems and opinion mining.
\begin{itemize}
\item \textbf{RQ1.} Which signal recovers hidden papers best: popularity, TF-IDF content, Personalized PageRank (PPR) or cocitation?
\item \textbf{RQ2.} Does a combination of signals beat the best single signal?
\end{itemize}

\section{Datasets}
\begin{itemize}
\item \textbf{OpenAlex + PeerRead}: 3,595 AI and ML papers with 13,481 citation edges; 75.80\% have an abstract.
\item \textbf{PeerRead} (ICLR 2017): 1,320 scored reviews, split 1,082 / 123 / 115.
\item \textbf{SciFact} (BEIR): 5,183 abstracts and 300 queries, to check search.
\end{itemize}

\section{Current experiment results}
We hide citations and ask each method to recover them. E2: cut at 2020, hide one third of the references of 219 later papers. E2b: predict what 250 authors cite for the first time after 2020. E3: random holdout of 442 edges. Metric: nDCG@10.

\begin{center}\small
\begin{tabular}{@{}lrrr@{}}\toprule
Method & E2 & E2b & E3\\\midrule
Content (TF-IDF) & 0.110 & 0.068 & 0.070\\
PPR & 0.121 & 0.047 & 0.174\\
Hybrid (0.7 content + 0.3 PPR) & 0.143 & 0.070 & 0.112\\
Cocitation & 0.173 & 0.099 & \textbf{0.305}\\
Fusion (RRF of cocitation and hybrid) & 0.207 & 0.113 & 0.259\\
Decision model (gradient boosting) & 0.224 & 0.132 & 0.297\\
Decision model + topic sentence & \textbf{0.237} & \textbf{0.182} & 0.304\\\bottomrule
\end{tabular}
\end{center}

\begin{itemize}
\item \textbf{RQ1:} cocitation is the best single signal on all three tasks.
\item \textbf{RQ2:} fusion beats the hybrid on all three tasks, but loses to cocitation alone on E3. A gradient boosting model trained on 2,081 tasks beats fusion on all three; with a topic sentence, all three gains have 95\% intervals above zero.
\item Language models used as rankers (Qwen2.5, GLM 5.3 Flash) did not beat the decision model and are more than 1,000 times slower.
\item Reviews: TF-IDF with logistic regression reaches macro F1 0.469 (majority class 0.181). Search: BM25 nDCG@10 0.651 on SciFact.
\end{itemize}

\section{System: current stage}
\textbf{Requirements.} Search and save papers, recommend with evidence, show the citation graph and reviews, organize projects, read PDFs with notes. Runs on one CPU machine.

\textbf{Design and implementation.} FastAPI backend, HTML and JavaScript interface, SQLite for libraries and notes. TF-IDF, BM25 and the citation graph are built in memory; NetworkX computes PageRank and communities. All views are working, and 574 automated tests pass.

\section{Next steps for the final report and presentation}
\begin{enumerate}
\item Repeat the time cutoff at other years to settle fusion versus cocitation.
\item Collect relevance labels from independent readers for project suggestions.
\item Finish the final report, the slides and the demo.
\end{enumerate}
\end{document}
